% ===========================================================================
%  Chapter 01 homework — Groups, fields and vector spaces
%  Garling, A Course in Galois Theory
%
%  DIVISION OF LABOUR (see AI_INSTRUCTIONS.md §7):
%    The assistant transcribes each assigned problem statement faithfully and
%    emits an empty, marked solution region beneath it.
%    The user types the mathematics inside those regions. No assistant writes
%    inside a SOLUTION region in normal mode.
% ===========================================================================
\documentclass[11pt]{article}
\usepackage{coursemacros}

\title{Chapter 01 Homework --- Groups, fields and vector spaces}
\author{Mikey Ferguson}
\date{\today}

\begin{document}
\maketitle

% ---------------------------------------------------------------------------
% Assistant: replace this block with one problem/solution pair per assigned
% exercise, in the order assigned, following the pattern below exactly.
% ---------------------------------------------------------------------------

\begin{problem}{01.3}
  Suppose that $H$ is a subgroup of $G$ of index $2$. Show that $H$ is a normal
  subgroup of $G$.
\end{problem}
% ===== SOLUTION 01.3 =====
We must show that $g h g^{-1} \in H$ for every $g \in G$ and every $h \in H$.
Let $h \in H$.

\medskip
\noindent\textbf{Case 1: $g \in H$.}
Then $g h g^{-1} \in H$ by closure; no normality is needed to guarantee this.

\medskip
\noindent\textbf{Case 2: $g \notin H$.}
Since $H$ has index $2$, there are exactly two cosets of $H$ in $G$, namely $H$ and
one other; let $gH$ be that disjoint coset of $H$, so $g \notin H$. Denote
$x = g h g^{-1}$.

Suppose $x \in gH$. Then there exists $h_1 \in H$ such that
\[
  g h g^{-1} \;=\; g h_1
  \iff g^{-1} \;=\; h^{-1} h_1
  \iff g \in H,
\]
a contradiction.

So since $x \notin gH$ due to the contradiction, $x \in H$, since $H$ and $gH$
partition $G$. So $x = g h g^{-1} \in H$ and $H$ is normal. \qed
% ===== END SOLUTION 01.3 =====

\end{document}
